\bookchapter{บทที่ 13: แนะนำแอปประเภทเครื่องมือ}

\chapterintro{บทนี้คัดสรรเครื่องมือด้านประสิทธิภาพ สุขภาพ และการเดินทาง พร้อมเสริมความสามารถของระบบและการทำงานข้ามอุปกรณ์: การทำความสะอาดและป้องกันสายรบกวนอย่าง Qingzhuo สวิตช์บล็อกที่มือถือมีในตัว การเตือนสภาพอากาศและแผ่นดินไหว การควบคุมคอมพิวเตอร์ระยะไกลอย่าง Sunlogin/ToDesk การทำงานออฟฟิศด้วย AI อย่าง Trae วอลเปเปอร์เคลื่อนไหว การส่งข้อมูลเร็วบน LAN อย่าง LocalSend คลิปบอร์ดข้ามอุปกรณ์ รหัสผ่านแบบไดนามิกอย่าง Aegis/Microsoft Authenticator ตารางงาน TickTick และการสแกนกระดาษ}

\booksection{เครื่องมือเพิ่มประสิทธิภาพ: จัดการเวลาและสมาธิของคุณ}

เครื่องมือเพิ่มประสิทธิภาพคือ “โครงกระดูก” ของชีวิตดิจิทัล ช่วยจัดระเบียบความคิด งาน และข้อมูลที่กระจัดกระจายให้เป็นระเบียบ ด้านล่างแนะนำตามสถานการณ์

\booksubsection{รายการสิ่งที่ต้องทำ (การจัดการงาน)}

\begin{itemize}
  \item TickTick เหมาะกับ: คนที่ต้องการบันทึกงาน ชีวิต การเรียน และตารางนัดหมายไว้ในที่เดียว เป็นหนึ่งในตัวแนะนำหลักของส่วนนี้ จุดเด่น: ซิงก์ข้ามมือถือ/คอมพิวเตอร์/แท็บเล็ต/Web รวมรายการ + ปฏิทิน + นาฬิกาโพโมโดโรไว้ในหนึ่งเดียว ป้อนภาษาธรรมชาติเช่น "ประชุมพรุ่งนี้บ่าย 3 โมง" "ส่งรายงานประจำสัปดาห์ทุกวันจันทร์" ระบบจะแยกแยะเวลาและกฎทำซ้ำอัตโนมัติ รองรับงานย่อย แท็ก ลำดับความสำคัญ ไฟล์แนบ และรายการที่ทำงานร่วมกัน เหมาะกับการจัดการ "ตารางงาน + ทำเครื่องหมายงาน + ทบทวนประจำวัน" แบบครบวงจร
  \item ใช้ TickTick ให้คล่องขึ้น: กล่องรับ: โยนทุกอย่างเข้า "วันนี้/กล่องรับ" ก่อน ตอนเย็นค่อยจัดเข้ารายการงาน ครอบครัว การเรียน
  \item มุมมองปฏิทิน: ลากการประชุมและกำหนดส่งลงปฏิทินเพื่อดูช่วงเวลาที่ขัดแย้ง เห็นภาพชัดกว่าดูแค่รายการสิ่งที่ต้องทำ
  \item งานซ้ำ: ตั้งกินยา จ่ายบิล ออกกำลังกายเป็นงานตามกฎ ไม่ต้องกรอกเองทุกครั้ง
  \item นาฬิกาโพโมโดโร: เปิดโหมดโฟกัสจากงานได้โดยตรง นับสถิติ "สัปดาห์นี้ทำอะไรเสร็จจริงบ้าง"
  \item เชื่อมโยงกับฟังก์ชันมือถือ: ผู้ช่วยเสียงพูด "เตือนฉัน..." จะเข้าการเตือนของระบบ กำหนดส่งสำคัญเปิดปฏิทินระบบคู่กันเป็นประกันสองชั้น
\end{itemize}

Todoist เหมาะกับ: ผู้ใช้ที่ชอบความเรียบง่ายและหลายแพลตฟอร์ม จุดเด่น: อินเทอร์เฟซเรียบสะอาด ลำดับชั้นโปรเจกต์ชัดเจน (งานงานย่อยงานย่อยย่อย) คะแนน karma จูงใจการสร้างนิสัย

Microsoft To Do เหมาะกับ: คนทำงานออฟฟิศ ผู้ใช้ Outlook จุดเด่น: ฟรีทั้งหมด ไม่มีโฆษณา ฟังก์ชัน "วันของฉัน" ล้างทุกวัน ลดภาระทางใจ ผสานกับระบบนิเวศ Microsoft อย่างลึกซึ้ง

โน้ต/การเตือนความจำของระบบ เหมาะกับ: ผู้ใช้เบา ๆ ที่ไม่ต้องการติดตั้งแอปบุคคลที่สาม จุดเด่น: ประสบการณ์ดั้งเดิม ไม่หน่วง เพิ่มด้วยเสียง Siri/Xiao Ai ได้โดยตรง เชื่อมโยงกับวิดเจ็ตของระบบ

\booksubsection{นาฬิกาโพโมโดโร (ฝึกสมาธิ)}

\begin{itemize}
  \item Forest เหมาะกับ: คนที่ห้ามใจไม่ไหว เล่นมือถือง่าย จุดเด่น: ปลูกต้นไม้ขณะโฟกัส ออกจากแอปก่อนเวลาต้นไม้จะเหี่ยวตาย ใช้ "ชีวิตเสมือน" สร้างการตอบรับเชิงบวก ปลูกร่วมกับเพื่อนได้
  \item Tide เหมาะกับ: คนที่ต้องการเสียงสีขาวช่วยโฟกัส/หลับ จุดเด่น: เสียงธรรมชาติ (ฝน คลื่น ป่า) + นาฬิกาโพโมโดโร + การทำสมาธิ รวมสามในหนึ่ง อินเทอร์เฟซเรียบง่ายเยียวยา
  \item หลักการ: โฟกัส 25 นาที + พัก 5 นาทีแบบคลาสสิก ปรับเป็น 50/10 หรือ 90/20 ตามสภาพร่างกายได้
\end{itemize}

\booksubsection{ปฏิทิน (เห็นภาพเวลา)}

\begin{itemize}
  \item ปฏิทินระบบ เหมาะกับ: ผู้ใช้ทุกคน จุดเด่น: พุชดั้งเดิมเสถียรที่สุด กินแบตน้อยมาก เพิ่มด้วยเสียง Siri/Xiao Ai ได้
  \item Google Calendar เหมาะกับ: ผู้ใช้หลายอุปกรณ์ คนที่ต้องแชร์ตารางกับผู้อื่น จุดเด่น: ฟังก์ชัน "เป้าหมาย" จัดเวลาออกกำลังกาย/เรียนให้อัตโนมัติ ผสานกับ Gmail แยกแยะกิจกรรมในอีเมลอย่างชาญฉลาด
\end{itemize}

\booksubsection{แอปโน้ต (สะสมความรู้)}

\begin{itemize}
  \item Notion เหมาะกับ: คนที่ต้องการจัดการโน้ต งาน ฐานข้อมูล wiki ไว้ในที่เดียว นักศึกษา ฟรีแลนซ์ ทีมที่ทำงานร่วมกัน จุดเด่น: คลังความรู้ครบวงจร แก้ไขแบบบล็อก สร้างหน้าส่วนตัว รายการอ่าน กระดานโปรเจกต์ ระบบนิเวศเทมเพลตอุดมสมบูรณ์
  \item Obsidian เหมาะกับ: ผู้จัดการความรู้เชิงลึก นักวิจัย นักเขียน จุดเด่น: เก็บ Markdown ในเครื่อง ลิงก์สองทาง + แผนภาพความสัมพันธ์ ทำให้โน้ตของคุณกลายเป็น "สมองที่สอง" ข้อมูลเป็นของคุณทั้งหมด
  \item flomo เหมาะกับ: คนที่อยากบันทึกแรงบันดาลใจเป็นชิ้น ๆ แต่ไม่อยากเขียนยาว จุดเด่น: โน้ตแบบการ์ด บันทึกความคิดเหมือนโพสต์ Moments รองรับแท็ก \# จัดหมวดหมู่ ย้อนดูตามไทม์ไลน์
  \item Apple Notes เหมาะกับ: ผู้ใช้ iPhone/Mac ความต้องการบันทึกเบา ๆ จุดเด่น: เปิดเร็ว รองรับสแกนเอกสาร วาดมือ เช็กลิสต์ ซิงก์ iCloud ฟรีและผสานกับระบบอย่างลึกซึ้ง
  \item Evernote เหมาะกับ: ผู้ใช้ขั้นสูงที่ต้องตัดเก็บเว็บจำนวนมาก จัดเก็บเอกสาร จุดเด่น: ปลั๊กอินตัดเก็บเว็บทรงพลัง OCR อ่านข้อความในรูปภาพ แบรนด์เก่าแก่เสถียร แต่เวอร์ชันพรีเมียมต้องเสียเงิน
\end{itemize}

\booksubsection{เครื่องมือ PDF}

\begin{itemize}
  \item PDF Expert เหมาะกับ: ผู้ใช้ iPad/iPhone คนทำธุรกิจที่ต้องมาร์กอัปสัญญา นักศึกษา จุดเด่น: ตัวแก้ไข PDF ที่ได้รับการยอมรับว่าดีที่สุดบน iOS ไฮไลต์ลื่นไหล รองรับลายเซ็น รวม/แยก แบบฟอร์ม
  \item WPS เหมาะกับ: คนทำงานที่ต้องจัดการ Word/Excel/PPT/PDF พร้อมกัน จุดเด่น: ชุดสำนักงานครบวงจร รองรับคลาวด์จีน/ไฟล์ WeChat ดี เวอร์ชันฟรีเพียงพอสำหรับการใช้งานประจำวัน
  \item Adobe Acrobat เหมาะกับ: คนที่ต้องการมาตรฐานรูปแบบ PDF สูง (กฎหมาย ออกแบบ สิ่งพิมพ์) จุดเด่น: มาตรฐานทางการ แปลงรูปแบบ PDF/A ไม่ผิดพลาด การซิงก์ข้ามอุปกรณ์และอิเล็กทรอนิกส์ซิกเนเจอร์มืออาชีพ
\end{itemize}

\booksubsection{ประเภทสแกน (เปลี่ยนกระดาษเป็นดิจิทัล)}

สัญญากระดาษ ใบเสร็จ การบ้าน บันทึกไวท์บอร์ด สำเนาบัตรประชาชน... กล้องมือถือสามารถแทนสแกนเนอร์ส่วนใหญ่ได้แล้ว ขั้นตอนหลัก: ถ่ายให้ชัด ครอปและปรับมุมอัตโนมัติ ส่งออก PDF/รูปภาพ OCR เป็นตัวเลือก

\begin{itemize}
  \item CamScanner เหมาะกับ: คนทำธุรกิจและนักศึกษาที่ต้องสแกนสัญญา ใบเสร็จ บัตรประจำตัวบ่อย ๆ จุดเด่น: ตรวจขอบอัตโนมัติ ปรับมุมมอง OCR อ่านข้อความ สร้าง PDF หลายหน้าและซิงก์คลาวด์
  \item Microsoft Lens เหมาะกับ: นักศึกษา คนที่ต้องสแกนไวท์บอร์ด/หนังสือเป็น Word จุดเด่น: ฟรีไม่มีโฆษณา ส่งออกไป OneNote/Word ได้ในคลิกเดียว อ่านข้อความแม่นยำสูง
  \item การสแกนใน Apple Notes เหมาะกับ: ผู้ใช้ iPhone ความต้องการสแกนเบา ๆ จุดเด่น: ไม่ต้องติดตั้งแอป คลิกไอคอนกล้องในโน้ตเพื่อสแกน ตัดเงาอัตโนมัติ มาร์กอัปลายเซ็น ฟรีและใช้ดี
  \item Baimiao / CamScanner เว็บ / กล้องแบรนด์ "สแกนเอกสาร": กล้องระบบของ Huawei, Xiaomi, OPPO, vivo ส่วนใหญ่มีโหมด "เอกสาร/สแกน" ในตัว Android ยังใช้ Google Camera หรือสแกนของระบบได้ เหมาะกับการถ่ายสัญญาหนึ่งหน้าฉุกเฉิน
  \item เคล็ดลับถ่ายรูป: แสงสม่ำเสมอ มือถือขนานกับกระดาษ จัดมุมทั้งสี่มุม เอกสารหลายหน้าถ่ายตามลำดับแล้วรวม PDF เอกสารสำคัญระวังลายน้ำและที่เก็บ ใช้เสร็จล้างจาก "ลบทิ้งล่าสุด" ได้หมด
\end{itemize}

\booksection{สุขภาพและการออกกำลังกาย}

แอปสุขภาพคือ “ผู้ช่วยส่วนตัว” ที่ถูกมองข้ามมากที่สุดของมือถือ—สามารถติดตามการนอน อัตราการเต้นของหัวใจ ปริมาณการออกกำลังกายอย่างต่อเนื่อง สะท้อนสภาพระยะยาวได้ดีกว่าการไปโรงพยาบาลเป็นครั้งคราว

\booksubsection{ออกกำลังกายและฟิตเนส}

\begin{itemize}
  \item Keep เหมาะกับ: คนที่ออกกำลังกายที่บ้าน มือใหม่ต้นทุนเป็นศูนย์ คนที่อยากลดไขมันและกระชับหุ่น จุดเด่น: คอร์สฟรีมากมาย ครอบคลุมตั้งแต่ HIIT ถึงโยคะ AI สร้างแผนฝึกอัตโนมัติ ชุมชนคึกคัก
  \item Nike Training Club (NTC) เหมาะกับ: คนที่ต้องการคำแนะนำจากเทรนเนอร์มืออาชีพ ต้องการฝึกอย่างเป็นระบบ จุดเด่น: จาก Nike ทางการ คอร์สฟรีทั้งหมด การสาธิตท่าโดยเทรนเนอร์มืออาชีพ จัดหมวดหมู่ตามกล้ามเนื้อ/อุปกรณ์ชัดเจน
  \item Apple Fitness เหมาะกับ: ผู้ใช้ Apple Watch คนที่อยากให้ระบบวงแหวนจูงใจให้ขยับ จุดเด่น: วงแหวนกิจกรรมสามวงทุกวัน (ขยับ/ออกกำลังกาย/ยืน) เป้าหมายเห็นภาพ เชื่อมโยงข้อมูลอัตราการเต้นของหัวใจจากนาฬิกาอย่างไร้รอยต่อ
\end{itemize}

\booksubsection{ติดตามการนอน}

\begin{itemize}
  \item AutoSleep เหมาะกับ: ผู้ใช้ Apple Watch คนที่อยากรู้สัดส่วนหลับลึก/หลับตื้นแต่ละคืน จุดเด่น: ใส่นาฬิกานอนจะติดตามอัตโนมัติ ไม่ต้องกดเริ่มหยุดเอง ให้คะแนนคุณภาพการนอนและคำแนะนำก่อนนอน
  \item Sleep Cycle เหมาะกับ: คนที่ไม่ใส่นาฬิกานอน อยากใช้มือถือวัด จุดเด่น: วิเคราะห์การพลิกตัว/เสียงกรนผ่านไมโครโฟน ปลุกอัจฉริยะในช่วงหลับตื้น ตื่นไม่ทรมาน
  \item Mi Fitness เหมาะกับ: ผู้ใช้ Mi Band/นาฬิกา Xiaomi จุดเด่น: ฟรีและเชื่อมโยงกับ Mi Band อย่างลึกซึ้ง ติดตามอัตราการเต้นของหัวใจ/ออกซิเจนในเลือด/การนอนครบวงจร คุ้มค่า
\end{itemize}

\booksubsection{ทำสมาธิและผ่อนคลาย}

\begin{itemize}
  \item Calm เหมาะกับ: พนักงานออฟฟิศเครียด คนนอนไม่หลับ จุดเด่น: การนำทำสมาธิ "ความสงบประจำวัน" + เรื่องเล่าก่อนนอน (อ่านโดยคนดัง) อินเทอร์เฟซสวย เนื้อหาภาษาอังกฤษอุดมสมบูรณ์
  \item Tide เหมาะกับ: ผู้ใช้ภาษาจีนที่ต้องการเสียงสีขาวช่วยโฟกัส/หลับ จุดเด่น: ของดีจากจีน เสียงธรรมชาติ + คอร์สทำสมาธิ + นาฬิกาโพโมโดโร การนำภาษาจีนคุ้นเคยกว่า
  \item Xiao Sleep เหมาะกับ: คนที่ต้องการเสียงช่วยนอน บันทึกเสียงพูดละเมอ/กรน จุดเด่น: เสียงสีขาวมากมายซ้อนรวมได้ (ฝน+พัดลม+กองไฟ) ตั้งเวลาปิดก่อนนอน บันทึกเสียงพูดละเมอทั้งสนุกและมีประโยชน์
\end{itemize}

\booksubsection{สุขภาพผู้หญิง}

\begin{itemize}
  \item Flo เหมาะกับ: คนที่ต้องการบันทึกประจำเดือน เตรียมตั้งครรภ์ หรือใส่ใจสุขภาพผู้หญิง จุดเด่น: AI พยากรณ์รอบเดือน/วันตกไข่ครั้งต่อไปแม่นยำ ถามในชุมชนแบบไม่เปิดเผยตัวตนได้ ดีไซน์อินเทอร์เฟซประณีต
  \item Meiyou เหมาะกับ: ผู้หญิงในจีน คนที่อยากดูชุมชนและบทความภาษาจีน จุดเด่น: เนื้อหาท้องถิ่นอุดมสมบูรณ์ ให้ความรู้ตามช่วงประจำเดือน/เตรียมตั้งครรภ์/ตั้งครรภ์/เลี้ยงลูก
\end{itemize}

\booksubsection{เตือนกินยา}

\begin{itemize}
  \item Medication Assistant เหมาะกับ: ผู้ป่วยโรคเรื้อรังที่ต้องกินยา ลูกที่ต้องดูแลผู้ใหญ่กินยา จุดเด่น: สแกนบาร์โค้ดระบุยา สร้างรายการยาอัตโนมัติ แจ้งเตือนตรงเวลา บันทึกการกินยาแชร์ให้ครอบครัวได้
\end{itemize}

\booksubsection{ศูนย์รวมข้อมูล}

\begin{itemize}
  \item Apple Health/Google Fit เหมาะกับ: ผู้ใช้สมาร์ทโฟนทุกคน จุดเด่น: รวมข้อมูลจากมือถือ สายรัดข้อมือ และแอปสุขภาพบุคคลที่สาม (ก้าวเดิน อัตราการเต้นของหัวใจ การนอน ประจำเดือน ฯลฯ) เป็น "ศูนย์กลางเวชระเบียนสุขภาพ" ที่เป็นหนึ่งเดียว
\end{itemize}

\booksection{การเดินทางและท่องเที่ยว}

ออกไปข้างนอก แอปดี ๆ ไม่กี่ตัวช่วยให้คุณไม่ต้องเจออุปสรรคมากมาย ด้านล่างแนะนำตลอดกระบวนการเดินทาง

\booksubsection{เที่ยวบินและรถไฟความเร็วสูง}

\begin{itemize}
  \item Umetrip เหมาะกับ: คนทำธุรกิจที่เดินทางเครื่องบินบ่อย จุดเด่น: ข้อมูลการบินทางการ สถานะเที่ยวบิน/เปลี่ยนประตูขึ้นเครื่อง/เตือนความล่าช้าแม่นยำที่สุด ดูรุ่นเครื่องบินและแผนที่ที่นั่งได้
  \item 12306 เหมาะกับ: คนเดินทางรถไฟในจีนทุกคน จุดเด่น: ช่องทางซื้อตั๋วทางการที่เชื่อถือได้เพียงหนึ่งเดียว รองรับซื้อตั๋วรอคิว (แย่งตั๋วอัตโนมัติเมื่อมีคนคืน) ตั๋วนักเรียน/ตั๋วกลุ่มสะดวก
  \item Flight Master เหมาะกับ: คนที่อ่อนไหวต่อความตรงเวลาของเที่ยวบิน คนที่ต้องรับส่งสนามบิน จุดเด่น: พยากรณ์ความล่าช้าของเที่ยวบินแม่นยำ แชร์แผนเดินทางให้ครอบครัวดูตำแหน่งและเวลาถึงแบบเรียลไทม์
\end{itemize}

\booksubsection{แผนที่และการนำทาง}

\begin{itemize}
  \item Amap เหมาะกับ: ผู้ใช้ในเมืองจีน คนขับรถ/เรียกแท็กซี่ จุดเด่น: จราจรเรียลไทม์แม่นยำ เลี่ยงรถติดเก่ง เรียกแท็กซี่/คาร์พูลครบวงจร ตัวเลือกต่อรถสาธารณะหลากหลาย
  \item Baidu Maps เหมาะกับ: คนที่ต้องการข้อมูลสถานที่และการนำทางในอาคาร จุดเด่น: แผนที่พาโนรามา/นำทาง AR แผนที่ในอาคารสถานที่ท่องเที่ยว/ห้างสรรพสินค้าละเอียด แพ็กเสียงหลากหลาย
  \item Google Maps เหมาะกับ: คนเดินทางต่างประเทศ จุดเด่น: ครอบคลุมทั่วโลกมากที่สุด รีวิวท้องถิ่นจริง Street View และแผนที่ออฟไลน์เป็นตัวช่วยในต่างแดน
\end{itemize}

\booksubsection{แปลภาษา}

\begin{itemize}
  \item Google Translate เหมาะกับ: คนเดินทางต่างประเทศ คนเรียนภาษา จุดเด่น: "การสนทนา" แปลสองทางเรียลไทม์ "กล้อง" แปลคำจากภาพทันที เป็นของดีเวลาเดินทาง รองรับดาวน์โหลดแพ็กภาษาออฟไลน์
\end{itemize}

\booksubsection{แปลงอัตราแลกเปลี่ยน}

\begin{itemize}
  \item Minimalist Exchange เหมาะกับ: คนเดินทางต่างประเทศ ช้อปปิ้งข้ามพรมแดน จุดเด่น: อินเทอร์เฟซเรียบสะอาดมาก รองรับกว่า 160 สกุลเงิน อัปเดตเรียลไทม์ เปรียบเทียบหลายสกุลเงินพร้อมกัน กดค้างตัวเลขเพื่อแปลงเร็ว
\end{itemize}

\booksubsection{สภาพอากาศ}

\begin{itemize}
  \item สภาพอากาศระบบ / China Weather เหมาะกับ: ทุกคน ตัวเลือกแรกสำหรับดูพยากรณ์และเตือนภัยทางการ จุดเด่น: แอปสภาพอากาศของผู้ผลิตเชื่อมโยงกับช่องทางทางการเช่น "China Weather" เตือนพายุ ไต้ฝุ่น คลื่นความร้อน จะแจ้งเตือนผ่านแบนเนอร์/การแจ้งเตือน มีอำนาจสูง
  \item Caiyun Weather เหมาะกับ: คนที่ต้องการพยากรณ์ฝนแม่นยำระดับนาที จุดเด่น: "อีกสองชั่วโมงฝนจะตกไหม" แม่นยำนำในวงการ แผนที่ฝนแบบไดนามิกเห็นชัด แจ้งเตือนก่อนฝนตก
  \item Carrot Weather เหมาะกับ: ผู้ชื่นชอบสภาพอากาศที่สนุกและวิดเจ็ตหลากหลาย จุดเด่น: AI พยากรณ์อากาศด้วยเสียงถากถาง วิดเจ็ตปรับแต่งมากมาย แผนเรดาร์มืออาชีพ เวอร์ชันเสียเงินฟังก์ชันครบกว่า
  \item จุดสำคัญเรื่องเตือนภัย: แอปสภาพอากาศต้องอนุญาตการแจ้งเตือนจึงจะได้รับการเตือน การเตือนแบ่งเป็นสี่ระดับ ฟ้า เหลือง ส้ม แดง สีแดงหมายถึงความเสี่ยงสภาพอากาศภัยพิบัติสูงสุด ควรปรับแผนเดินทางตามคำเตือน วิธีเปิดใช้งานเพิ่มเติมดูส่วนที่ 7 ของบทนี้ "เตือนสภาพอากาศ / เตือนภัยพิบัติธรรมชาติ / เตือนแผ่นดินไหว"
\end{itemize}

\booksubsection{คู่มือท่องเที่ยว}

\begin{itemize}
  \item Mafengwo เหมาะกับ: คนที่ชอบทำคู่มือละเอียด เดินทางด้วยรถส่วนตัว/ท่องเที่ยวเชิงลึก จุดเด่น: เนื้อหาบันทึกการเดินทางจากผู้ใช้ดั้งเดิมอุดมสมบูรณ์ วิธีเที่ยว/เส้นทาง/คู่มือเลี่ยงหลุมบ่อละเอียด ดาวน์โหลดคู่มือออฟไลน์ได้
  \item Xiaohongshu เหมาะกับ: นักเดินทางรุ่นใหม่ที่อยากดูประสบการณ์จริง หาจุดเช็กอิน/ของกินเฉพาะกลุ่ม จุดเด่น: รูปภาพ ข้อความ วิดีโอสั้นดูเข้าใจง่าย เนื้อหาเช่น "คู่มือ 3 วันในเมือง XX" อัปเดตเร็ว ทันสมัย ค้นหาดี
\end{itemize}

\booksection{ทำความรู้จักเครื่องมืออัตโนมัติ}

ระบบอัตโนมัติคือคานงัดที่ “ตั้งค่าครั้งเดียว ประหยัดระยะยาว” สามารถบีบอัดการกระทำที่ทำซ้ำทุกวันหลายสิบครั้งให้เหลือศูนย์ ส่วนนี้แนะนำแนวคิดเท่านั้น วิธีใช้และบทเรียนละเอียดดู บทที่ 23: ความคิดเชิงโปรแกรมและการเริ่มต้นระบบอัตโนมัติ

\begin{itemize}
  \item iOS "คำสั่งลัด": เครื่องมืออัตโนมัติทางการจาก Apple ไม่ต้องเขียนโค้ด ทำชุดการทำงานเช่น "ถึงบ้านเปิด WiFi + หรี่แสง + พยากรณ์อากาศ" ได้ในคลิกเดียว
  \item ระบบอัตโนมัติบน Android: การเชื่อมโยงสถานการณ์ระดับระบบ (เช่น "งานอัตโนมัติ" ของ Xiaomi, "สถานการณ์อัจฉริยะ" ของ Huawei) เริ่มต้นง่ายที่สุด Tasker ทรงพลังแต่เส้นโค้งการเรียนรู้ชัน เหมาะกับคนชอบเทคนิค
  \item แนวคิดหลัก: หา "การกระทำที่ต้องทำซ้ำทุกวัน" ก่อน แล้วคิดว่า "ใช้เงื่อนไขทริกเกอร์เดียว + หลายการกระทำแทนได้ไหม" ตัวอย่างละเอียดและวิธีขั้นสูงดูบทที่ 23
\end{itemize}

\booksection{เบราว์เซอร์และเครื่องมือค้นหา: ประตูสู่ข้อมูลของคุณ}

เบราว์เซอร์คือหน้าต่างสู่อินเทอร์เน็ต เครื่องมือค้นหาคือเข็มทิศนำทางในทะเลข้อมูล ใช้ให้เป็นจะช่วยให้คุณไม่ต้องหลงทาง

\booksubsection{หลักการทำงานของเบราว์เซอร์:}

เมื่อคุณพิมพ์ที่อยู่เว็บในเบราว์เซอร์ เบราว์เซอร์จะ: (1)สอบถามที่อยู่ IP ของโดเมนจากเซิร์ฟเวอร์ DNS (2)ส่งคำขอ HTTP/HTTPS ไปยังเซิร์ฟเวอร์ที่อยู่ IP นั้น (3)เซิร์ฟเวอร์ส่งโค้ดเว็บเพจกลับมา (HTML, CSS, JavaScript) (4)เบราว์เซอร์แปลงโค้ดและแสดงผลเป็นหน้าที่คุณเห็น

\booksubsection{เบราว์เซอร์หลักที่แนะนำ:}

\begin{itemize}
  \item Microsoft Edge: ใช้ Chromium เป็นแกน เร็ว เข้ากันได้ดี มีผู้ช่วย Copilot AI ในตัว สรุปเว็บเพจ แปลภาษา สนทนาได้ รองรับการซิงก์ข้ามอุปกรณ์ ผู้เขียนแนะนำเป็นเบราว์เซอร์หลักอย่างยิ่ง
  \item Safari: เบราว์เซอร์เริ่มต้นของอุปกรณ์ Apple ผสานกับระบบอย่างลึกซึ้ง ประหยัดแบต มีประสิทธิภาพ ปกป้องความเป็นส่วนตัวดี แต่ปลั๊กอินเสริมน้อย
  \item Chrome: ใช้มากที่สุดในโลก ส่วนเสริมหลากหลาย แต่กินหน่วยความจำมาก และค่อนข้างเปลืองแบตบนมือถือ
\end{itemize}

\booksubsection{การเลือกเครื่องมือค้นหา:}

\begin{itemize}
  \item Bing: เครื่องมือค้นหาของ Microsoft ผสานกับ Edge อย่างลึกซึ้ง ฟังก์ชันสรุป AI ให้คำตอบตรง ๆ ไม่ใช่แค่ลิงก์ ค้นหาได้ประสิทธิภาพมาก ประสบการณ์ค้นหารูปภาพก็ดี
  \item Yandex: เครื่องมือค้นหาใหญ่ที่สุดของรัสเซีย ค้นหารูปภาพและค้นภาพย้อนกลับเก่งมาก หาภาพคล้าย ระบุวัตถุได้ดีมาก
  \item Baidu: กำลังหลักในการค้นข้อมูลในจีน แต่โฆษณามาก ต้องใช้วิจารณญาณ
  \item Google: เครื่องมือค้นหาแม่นยำที่สุดในโลก แต่ในจีนต้องใช้เครือข่ายพิเศษจึงจะเข้าถึงได้
\end{itemize}

\booksubsection{คู่มือค้นข้อมูลอย่างมีประสิทธิภาพ:}

\begin{itemize}
  \item ใช้ AI ค้นหาให้เป็น: สรุป AI ของ Bing, Copilot ของ Edge ตอบคำถามได้ตรง ๆ ช่วยคัดสาระสำคัญจากผลลัพธ์มากมาย คล่องกว่าคลิกลิงก์ทีละอันมาก
  \item แยกแยะข้อมูลเท็จและโฆษณา: รายการในผลค้นหาที่ติดป้าย "โฆษณา" หรือ "ผู้สนับสนุน" เป็นการโปรโมตเสียเงิน ไม่ใช่ผลค้นหาธรรมชาติ ใส่ใจแหล่งข้อมูลที่น่าเชื่อถือ เลือกเว็บทางการและสื่อที่มีชื่อเสียงก่อน
  \item ดาวน์โหลดแอปให้ถูกวิธี: พิมพ์ "ชื่อแอป + เว็บทางการ" ในเครื่องมือค้นหา เช่น "เว็บทางการ WeChat" คลิกผลที่มีโดเมนทางการก่อน อย่าติดตั้งซอฟต์แวร์จากเว็บดาวน์โหลดบุคคลที่สาม เพราะมักแถมซอฟต์แวร์รบกวน
  \item เคล็ดลับแยกแยะเว็บทางการ: ดาวน์โหลด Windows 11 ควรเข้า microsoft.com ดาวน์โหลด Chrome ควรเข้า google.com ดูว่าโดเมนตรงกับแบรนด์ไหม เป็นวิธีแยกแยะพื้นฐานที่สุด
\end{itemize}

\booksubsection{ความรู้เรื่องโดเมนเว็บไซต์:}

โดเมนคือ “ป้ายที่อยู่” ของเว็บไซต์ อ่านออกจะช่วยให้คุณตัดสินประเภทและแหล่งที่มาของเว็บได้:

\begin{itemize}
  \item ความหมายของนามสกุลทั่วไป: .com — องค์กรธุรกิจ พบบ่อยที่สุด
  \item .cn — โดเมนระดับประเทศจีน
  \item .org — องค์กรไม่แสวงผลกำไร
  \item .net — องค์กรบริการเครือข่าย
  \item .edu — สถาบันการศึกษา
  \item .gov — หน่วยงานรัฐบาล
  \item .xyz/.tech/.info — โดเมนทั่วไปน้องใหม่ เกณฑ์ต่ำ ต้องระวัง
\end{itemize}

โดเมนประเทศ/ภูมิภาค: .jp — ญี่ปุ่น

.kr — เกาหลี

.uk — อังกฤษ

.de — เยอรมนี

.tw — ไต้หวัน

.hk — ฮ่องกง

ข้อควรระวัง: เว็บที่มีนามสกุล .gov และ .edu มักน่าเชื่อถือกว่า เพราะการสมัครเข้มงวด โดเมนราคาถูกอย่าง .xyz, .top ถูกเว็บฟิชชิงนำไปใช้ง่าย เข้าชมต้องระวังเป็นพิเศษ

\booksubsection{การบล็อกโฆษณา:}

\begin{itemize}
  \item AdGuard เหมาะกับ: ผู้ใช้ทุกคนที่รำคาญป๊อปอัป โฆษณาก่อนวิดีโอ โฆษณาในแอป จุดเด่น: บล็อกโฆษณาในระดับระบบ (ไม่จำกัดเฉพาะเบราว์เซอร์) กรองแบนเนอร์ในแอปและการติดตามหน้าเปิด ประหยัดข้อมูล โหลดหน้าเร็วขึ้น เวอร์ชันเสียเงินรองรับทุกอุปกรณ์
\end{itemize}

\booksubsection{เบราว์เซอร์ปกป้องความเป็นส่วนตัว:}

\begin{itemize}
  \item Brave เหมาะกับ: ผู้ใช้ทั่วไปที่ใส่ใจความเป็นส่วนตัว ไม่อยากถูกติดตามด้วยโฆษณา จุดเด่น: ใช้ Chromium เป็นแกน บล็อกโฆษณาและตัวติดตามตั้งแต่แกะกล่อง มีโหมดไม่ระบุตัวตน Tor ในตัว เร็วกว่า Chrome และใช้หน่วยความจำน้อยกว่า
  \item DuckDuckGo เหมาะกับ: ผู้ใช้ที่อยาก "ค้นคำเดียวแล้วไม่ถูกโฆษณาตามทั่วเน็ต" จุดเด่น: มีเครื่องมือค้นหาชื่อเดียวกันในตัว (ไม่บันทึกประวัติค้นหา) ล้างร่องรอยการท่องเว็บทั้งหมดในคลิกเดียว ("ปุ่มเผา") ประสบการณ์ข้ามแพลตฟอร์มเป็นหนึ่งเดียว
\end{itemize}

\booksubsection{ตารางเปรียบเทียบเบราว์เซอร์:}

\begin{center}
\begin{tabularx}{\linewidth}{|X|X|X|}
\hline
\textbf{เบราว์เซอร์} & \textbf{จุดเด่น} & \textbf{เหมาะกับใคร} \tabularnewline
\hline
Microsoft Edge & แกน Chromium + Copilot AI ในตัว สรุปเว็บ/แปล/สนทนาได้ & คนทำงานที่เน้นประสิทธิภาพ ผู้ใช้ AI หนัก ๆ (แนะนำเป็นตัวหลัก) \tabularnewline
Safari & ผสานกับระบบนิเวศ Apple อย่างลึกซึ้ง ประหยัดแบต มีประสิทธิภาพ ปกป้องความเป็นส่วนตัวแรง & ผู้ใช้ iPhone/Mac ที่ให้ความสำคัญกับอายุแบตและการเชื่อมโยงระบบ \tabularnewline
Chrome & ใช้มากที่สุดในโลก ระบบนิเวศปลั๊กอินอุดมสมบูรณ์ที่สุด & ผู้ใช้ขั้นสูงที่ต้องการปลั๊กอินจำนวนมาก คนที่ซิงก์ข้ามอุปกรณ์ \tabularnewline
Brave & บล็อกโฆษณาและตัวติดตามตั้งแต่แกะกล่อง เร็วกว่า & ผู้ใช้ที่ใส่ใจความเป็นส่วนตัว อยากบล็อกโฆษณาโดยไม่ติดปลั๊กอิน \tabularnewline
DuckDuckGo & เครื่องมือค้นหาชื่อเดียวกันไม่บันทึกประวัติ ล้างร่องรอยในคลิกเดียว & ผู้ใช้ที่อ่อนไหวต่อความเป็นส่วนตัว ไม่อยากถูกติดตามแบบส่วนบุคคล \tabularnewline
\hline
\end{tabularx}
\end{center}

\begin{tcbpractice}{ลองสร้างขั้นตอนอัตโนมัติ "ผ่อนคลายก่อนนอน"}
\begin{enumerate}
  \item สร้างเงื่อนไขทริกเกอร์: เวลาที่กำหนด (เช่น 22:30)
  \item เพิ่มการกระทำ: หรี่เสียง > เปิดโหมดห้ามรบกวน > เล่นเพลย์ลิสต์ผ่อนคลาย > หรี่แสง
  \item บันทึกและทดสอบ
\end{enumerate}

\end{tcbpractice}

\booksection{เครื่องมือระบบใช้งานจริง: ทำความสะอาด ป้องกันสายรบกวน และการป้องกันความปลอดภัย}

นอกจากแอปประสิทธิภาพและสุขภาพ สิ่งที่มีผลต่อประสบการณ์ประจำวันจริง ๆ มักเป็นสามเรื่อง: “ถูกรบกวนน้อย พื้นที่เพียงพอ มิจฉาชีพเข้าไม่ถึง” ส่วนนี้จะสำรวจแอปใช้งานจริงทั่วไปก่อน แล้วขยายเรื่องฟังก์ชันบล็อกสายรบกวนที่มือถือมีในตัวโดยเฉพาะ—หลายคนซื้อมือถือมาหลายปี แต่ไม่เคยเปิดสวิตช์เหล่านี้เลย

\booksubsection{ทำความสะอาดระบบและเพิ่มพื้นที่จัดเก็บ}

\begin{itemize}
  \item Qingzhuo เหมาะกับ: ผู้ใช้ Android ที่อยากทำความสะอาดแคช ไฟล์ที่เหลือ โฟลเดอร์ว่างเบา ๆ แต่ไม่อยากถูก "Mobile Manager" ชุดใหญ่รบกวน จุดเด่น: ขนาดเล็ก ไม่มีโฆษณารบกวน รองรับสแกนเชิงลึกหาไฟล์ที่เหลือของแอป หาไฟล์ใหญ่ วิเคราะห์การใช้พื้นที่ ประหยัดกว่า "เร่งความเร็วในคลิกเดียว" เชิงพาณิชย์ เหมาะกับผู้ใช้ที่ทำเองเก่งสำหรับการบำรุงรักษาตามปกติ หมายเหตุ: ก่อนทำความสะอาดให้ตรวจสอบเส้นทางก่อน อย่าติ๊กเลือกทั้งหมดในโฟลเดอร์ WeChat แกลเลอรี เอกสาร แคชระดับระบบควรใช้ตัวทำความสะอาดของแบรนด์ก่อน เครื่องมือบุคคลที่สามเป็นเพียงตัวเสริม
  \item ตัวทำความสะอาดพื้นที่ของระบบ เหมาะกับ: ทุกคน โดยเฉพาะผู้ใหญ่และผู้ใช้ที่ไม่อยากยุ่งยาก จุดเด่น: Android [การตั้งค่า] > [พื้นที่จัดเก็บ] > [ทำความสะอาดและเร่งความเร็ว] iOS [การตั้งค่า] > [ทั่วไป] > [พื้นที่จัดเก็บ iPhone] ระบบเข้าใจโครงสร้างแคชของแอปตัวเองดีกว่า ความเสี่ยงลบผิดต่ำสุด ควรเป็นตัวเลือกแรก
  \item ตัวจัดการพื้นที่ในตัวของ WeChat/QQ เหมาะกับ: คนแชทเยอะ จุดเด่น: WeChat [ฉัน] > [การตั้งค่า] > [ทั่วไป] > [พื้นที่จัดเก็บ] ล้างละเอียดตามแชทและประเภทไฟล์ได้ ปลอดภัยกว่าเครื่องมือ "ล้าง WeChat" ภายนอก
  \item หลักการ: ใช้ของระบบ/แอปในตัวได้ก็อย่าใช้บุคคลที่สาม เลือกตามเส้นทาง/ประเภทได้ก็อย่า "ล้างทั้งหมดในคลิกเดียว" เป้าหมายของการทำความสะอาดคือคืนพื้นที่ ไม่ใช่กวาดข้อมูลสำคัญไปด้วย
\end{itemize}

\booksubsection{แอปป้องกันสายรบกวนและความปลอดภัย}

\begin{itemize}
  \item National Anti-Fraud Center เหมาะกับ: ทุกคน โดยเฉพาะผู้ใหญ่ ตำแหน่งที่เกี่ยวกับการเงิน คนที่ช้อปออนไลน์บ่อย/รับสายแปลกหน้า จุดเด่น: ทางการจัดทำ เปิดเตือนสายเข้า เตือนความเสี่ยงข้อความ ตรวจเว็บน่าสงสัย เชื่อมโยงกับฐานข้อมูลป้องกันการทุจริตของตำรวจ เป็นหนึ่งใน "สามชั้นป้องกัน" ที่ได้รับการรับรองจากทางการมากที่สุด
  \item ศูนย์ความปลอดภัยของแบรนด์ / Mobile Manager เหมาะกับ: ผู้ใช้แบรนด์นั้น จุดเด่น: Huawei Mobile Manager, Xiaomi Security Center, OPPO Mobile Manager, vivo Mobile Manager ฯลฯ รวมบล็อกสายรบกวน สแกนไวรัส จัดการสิทธิ์ ป้องกันการชำระเงิน มักฟรีและทำงานร่วมกับระบบล่างได้ดีกว่า ไม่แนะนำติดตั้งชุดความปลอดภัยเชิงพาณิชย์ชุดที่สองซึ่งอาจทำให้สิทธิ์ขัดแย้งกัน
  \item บริการป้องกันสายรบกวนของผู้ให้บริการเครือข่าย เหมาะกับ: ผู้ใช้ที่ถูกโทรออกถี่ ๆ หรือสายเสียง AI ถล่ม จุดเด่น: China Mobile ส่ง KTFSR ไป 10086 เพื่อเปิดใช้ China Telecom ติดตาม "Tianyi Anti-Harassment" China Unicom ตั้งค่าใน "Unicom Mobile Manager" การบล็อกเกิดที่ฝั่งเครือข่าย มือถือไม่ดังก็บล็อกสายรบกวนได้จำนวนมาก
  \item AdGuard เหมาะกับ: ผู้ใช้ทุกคนที่รำคาญป๊อปอัป โฆษณาก่อนวิดีโอ โฆษณาในแอป จุดเด่น: บล็อกโฆษณาในระดับระบบ (ไม่จำกัดเฉพาะเบราว์เซอร์) กรองแบนเนอร์ในแอปและการติดตามหน้าเปิด ประหยัดข้อมูล โหลดหน้าเร็วขึ้น เวอร์ชันเสียเงินรองรับทุกอุปกรณ์
  \item ระวังชุด "ทำความสะอาดเร่งความเร็ว / บล็อกสายรบกวน" จากบุคคลที่สาม: ตัวจัดการบางตัวทำงานเบื้องหลังตลอด ขอสิทธิ์รายชื่อและข้อความซ้ำ ๆ ความสามารถบล็อกอาจไม่ดีกว่าระบบ แถมเปลืองแบตและเสี่ยงความเป็นส่วนตัว เปิดฟังก์ชันในตัวของระบบก่อน แล้วค่อยพิจารณาแอปป้องกันการทุจริตทางการ
\end{itemize}

\textbf{หลักการเลือกแอปใช้งานจริง}\par

\begin{itemize}
  \item ระบบมาก่อน: บล็อก ทำความสะอาด จัดการสิทธิ์ เปิดสวิตช์ในตัวของมือถือก่อน
  \item เสริมด้วยทางการ: National Anti-Fraud Center และบริการป้องกันสายรบกวนของเครือข่าย เป็นช่องทางทางการ ซ้อนได้
  \item บุคคลที่สามอย่างพอดี: ฟังก์ชันประเภทเดียวกันติดตั้งตัวเดียวก็พอ หลีกเลี่ยงสิทธิ์ขัดแย้งและการแย่งทรัพยากรเบื้องหลัง
  \item น้อยคือมาก: ยิ่งติดตั้งเครื่องมือมาก การแจ้งเตือนและสิทธิ์ยิ่งกว้าง โอกาสที่ "ตัวเครื่องมือเอง" จะรบกวนยิ่งสูง
\end{itemize}

\booksubsection{ฟังก์ชันบล็อกสายรบกวนในตัวของมือถือ: แต่ละแบรนด์เปิดอย่างไร}

ด้านล่างแสดงรายการทางเข้าตามแบรนด์ ชื่ออาจต่างกันเล็กน้อยตามเวอร์ชันระบบ ค้นหา “สายรบกวน” “บล็อก” “สายเข้าแปลกหน้า” ในการตั้งค่าได้โดยตรง เสริมกับบทที่ 8 “การจัดการสายรบกวน”: ที่นี่เน้นตำแหน่งสวิตช์ ที่นั่นเน้นการติดป้ายเบอร์และการอุทธรณ์ป้ายผิด

\begin{center}
\begin{tabularx}{\linewidth}{|X|X|X|}
\hline
\textbf{แพลตฟอร์ม} & \textbf{เส้นทางเข้า} & \textbf{รายการที่แนะนำให้เปิด} \tabularnewline
\hline
iPhone / iOS & {}[การตั้งค่า] > [โทรศัพท์] / [ข้อความ] & ปิดเสียงสายจากคนแปลกหน้า รายชื่อที่ถูกบล็อก "กรองผู้ส่งที่ไม่รู้จัก" ในข้อความ \tabularnewline
Xiaomi / Redmi & แอปโทรศัพท์ > การตั้งค่า > บล็อกสายรบกวน หรือ ศูนย์ความปลอดภัย & บล็อกอัจฉริยะ เบอร์แปลกหน้าดังครั้งเดียวแล้วตัด กรองคำสำคัญในข้อความ \tabularnewline
Huawei / Honor & โทรศัพท์ > บล็อกสายรบกวน Mobile Manager > บล็อกสายรบกวน & ระบุสายเข้า บัญชีดำ บล็อกข้อความอัจฉริยะ \tabularnewline
OPPO / realme / OnePlus & แอปโทรศัพท์ > บล็อกสายรบกวน หรือ เครื่องนี้ > บล็อกสายรบกวน & บล็อกสายเข้า ระบุสายหลอกลวง กฎบล็อกข้อความ \tabularnewline
vivo / iQOO & โทรศัพท์ > บล็อกสายรบกวน i Manager & บล็อกสายรบกวน เตือนสายหลอกลวง กรองข้อความ \tabularnewline
Samsung & โทรศัพท์ > การตั้งค่า > บล็อกหมายเลข / แสดงสายเข้าและสายขยะ & บล็อกหมายเลขที่ไม่รู้จัก ระบุสายขยะ (การรองรับต่างกันตามพื้นที่) \tabularnewline
\hline
\end{tabularx}
\end{center}

\begin{tcbinfo}{คำแนะนำการซ้อนสามชั้นป้องกัน}
แนะนำประจำวัน “บล็อกของระบบ + ป้องกันสายรบกวนของเครือข่าย + National Anti-Fraud Center” ซ้อนสามชั้น: ระบบจัดการป้ายและบัญชีดำ เครือข่ายบล็อกสายโทรออกถี่ ๆ ศูนย์ป้องกันการทุจริตเสริมเตือนทางการ ไม่ต้องเปิดทุกชั้นให้เข้มที่สุด เพราะอาจบล็อกสายพัสดุ อาหารเดลิเวอรี หมอผิด—อนุญาตให้ผ่านสำหรับรายชื่อติดต่อและผู้ติดต่อล่าสุดได้

\end{tcbinfo}

\begin{tcbpractice}{ลงมือทำ: เปิดบล็อกสายรบกวนวันนี้}
\begin{enumerate}
  \item เปิดการตั้งค่าแอปโทรศัพท์ หา "บล็อกสายรบกวน / ปิดเสียงสายแปลกหน้า" เปิดบล็อกอัจฉริยะ
  \item ตรวจมือถือของพ่อแม่ด้วย: เส้นทางเดียวกัน + เปิดบริการป้องกันสายรบกวนของเครือข่าย (China Mobile ส่ง KTFSR ไป 10086 ได้)
  \item ติดตั้ง "National Anti-Fraud Center" จากแอปสโตร์ ลงทะเบียนและเปิดเตือนสายเข้า
  \item ลองโทรไปเบอร์ทดสอบที่ถูกติดป้าย (หรือดูบันทึกการบล็อก) เพื่อยืนยันว่าฟังก์ชันทำงาน
\end{enumerate}

\end{tcbpractice}

\booksection{เตือนสภาพอากาศ เตือนภัยพิบัติธรรมชาติ และเตือนแผ่นดินไหว}

การบล็อกสายรบกวนป้องกัน “การหลอกลวงทางโทรศัพท์” ฟังก์ชันเตือนภัยป้องกัน “ภัยธรรมชาติและสภาพอากาศสุดขั้ว” หลายคนติดตั้งแอปสภาพอากาศแค่ดูอุณหภูมิ แต่ไม่เปิดการแจ้งเตือน ไม่เปิดเตือนแผ่นดินไหว—ก่อนพายุฝนสีแดง ไต้ฝุ่น แผ่นดินไหว ช่วงวินาทีถึงสิบกว่าวินาที มือถือช่วยให้คุณได้เวลาอันมีค่าในการหลบภัย

\booksubsection{เตือนสภาพอากาศ: รับการเตือนพายุ ไต้ฝุ่น คลื่นความร้อนอย่างไร}

\begin{itemize}
  \item ช่องทางทางการมาก่อน: การเตือนจาก China Meteorological Administration / สำนักอุตุนิยมวิทยากลาง / สำนักอุตุฯ แต่ละมณฑล มีอำนาจสูงสุด ติดตามบัญชีทางการ "China Weather" บน WeChat, Weibo หรือดาวน์โหลดแอปทางการเช่น "China Weather"
  \item แอปสภาพอากาศของระบบ: Huawei Weather, Xiaomi Weather, OPPO Weather, vivo Weather, Apple "Weather" ฯลฯ ส่วนใหญ่จะแจ้งเตือนจากสำนักอุตุฯ ท้องถิ่น เข้าการตั้งค่าแอปสภาพอากาศ ยืนยันว่าเปิดการแจ้งเตือน "เตือนสภาพอากาศ / สภาพอากาศภัยพิบัติ" แล้ว
  \item Caiyun Weather ฯลฯ พยากรณ์ฝนระดับนาที: เหมาะกับก่อนเดินทางไปทำงาน รับส่งลูก ตัดสินว่า "ครึ่งชั่วโมงนี้ต้องพกร่มไหม" ความสามารถเตือนภัยยังยึดสำนักอุตุฯ ทางการเป็นหลัก
  \item ความหมายสีการเตือน (ทั่วไป): สีฟ้า: ระวัง เดินทางตามปกติทั่วไป ติดตามอัปเดต
  \item สีเหลือง: เตรียมรับมือมากขึ้น ลดกิจกรรมกลางแจ้งที่ไม่จำเป็น
  \item สีส้ม: ร้ายแรง อยู่ในอาคารให้มากที่สุด อยู่ห่างจากพื้นที่อันตราย
  \item สีแดง: ร้ายแรงเป็นพิเศษ ทำตามคำสั่งหยุดเรียน หยุดงาน หลบภัย อย่าฝ่าน้ำท่วม น้ำป่า ถนนลมแรง
\end{itemize}

ประเภทการเตือนทั่วไป: พายุฝน หิมะตกหนัก ไต้ฝุ่น คลื่นความร้อน คลื่นความหนาว ลมแรง ฟ้าผ่า ลูกเห็บ ถนนเป็นน้ำแข็ง หมอกหนัก ความแห้งแล้ง ความเสี่ยงภัยทางธรณีวิทยา

\booksubsection{เตือนภัยพิบัติธรรมชาติ: ไม่ใช่แค่ “ฝนตก”}

\begin{itemize}
  \item ผู้ประกาศ: สำนักอุตุฯ รับผิดชอบด้านสภาพอากาศ ภัยทางธรณี น้ำป่า ฯลฯ มักร่วมประกาศโดยทรัพยากรธรรมชาติ ชลประทาน หน่วยฉุกเฉิน รัฐบาลท้องถิ่นก็ประกาศผ่านวิทยุฉุกเฉิน ข้อความ 12379 ฯลฯ
  \item รับหลายช่องทางอย่างไร: การแจ้งเตือนระบบมือถือ + เตือนภัยในแอปสภาพอากาศ
  \item WeChat ติดตามบัญชีทางการ "สำนักจัดการเหตุฉุกเฉิน" "สำนักอุตุฯ" ในพื้นที่
  \item ช่องทางแจ้งเตือนแห่งชาติ 12379 (ข้อความ/บริการ ตามที่เปิดใช้ในพื้นที่)
  \item ทีวี วิทยุฉุกเฉิน ลำโพงหมู่บ้าน/ชุมชน—บนดอย ช่วงน้ำหลาก อย่าจ้องแต่มือถือ
\end{itemize}

ข้อควรระวังกลางแจ้งและเดินทาง: วันไต้ฝุ่นอย่าอยู่ใกล้ป้ายโฆษณา อย่าเดินใต้ต้นไม้ ฝนตกหนักอย่าเข้าที่จอดรถใต้ดิน อุโมงค์ สะพานลอด ท่องเที่ยวบนดอยเมื่อได้รับเตือน "ความเสี่ยงภัยทางธรณี" ให้ออกจากไหล่เขา หุบเขาทันที

เชื่อมโยงกับคู่มือปฐมพยาบาล: หลังภัยพิบัติขาดการติดต่อ มือถือแบตหมด ต้องระบุตำแหน่งขอความช่วยเหลือ ดู บทที่ 32 สถานการณ์แผ่นดินไหว/ภัยพิบัติ และ บทที่ 7 การสื่อสารผ่านดาวเทียมและไม่มีเครือข่าย

\booksubsection{เตือนแผ่นดินไหว: ไม่เหมือน “พยากรณ์แผ่นดินไหว”}

การเตือนแผ่นดินไหว (EEW) คือการแจ้งเตือนล่วงหน้าหลายวินาทีถึงสิบกว่าวินาทีไปยังพื้นที่ที่อาจได้รับผลกระทบ ก่อนที่คลื่นแผ่นดินไหว (โดยเฉพาะคลื่น S ที่ทำลายล้างมากกว่า) จะถึงเป้าหมาย โดยใช้หลักที่คลื่นแม่เหล็กไฟฟ้าเร็วกว่าคลื่นแผ่นดินไหว และสถานีตรวจจับคลื่น P ได้เร็วกว่า มันไม่สามารถพยากรณ์ได้ว่าแผ่นดินไหวจะเกิดไหม แต่ช่วยยืดเวลา “หลบใกล้ตัว ตัดแหล่งไฟ วิ่งออกจากลิฟต์” ยิ่งใกล้ศูนย์กลาง ยิ่งมีเวลาเตือนน้อย อาจไม่ทันด้วยซ้ำ

\begin{itemize}
  \item ทางการและอุตสาหกรรม: ในจีนแผ่นดินใหญ่ ผู้ให้บริการที่พบบ่อย ได้แก่ ความสามารถของ China Earthquake Early Warning และบริการเตือนภัยที่ผลักดันโดยสถาบันเช่น Chengdu Institute of Disaster Prevention ผู้ผลิตบางรายติดตั้งล่วงหน้าหรือรับการเตือนผ่านเครือข่ายได้
  \item แอปบุคคลที่สาม: ค้นหา "เตือนแผ่นดินไหว" ในแอปสโตร์ (ก่อนติดตั้งดูผู้พัฒนาให้ดี อย่าหลงแอป "เตือนเท็จ" ที่ขอสิทธิ์ไม่เกี่ยวข้องมากมาย)
  \item เมื่อได้รับเตือนทำอย่างไร: ในอาคารให้หลบใกล้ตัว (ใต้โต๊ะ มุมผนังรับน้ำหนัก ข้างเฟอร์นิเจอร์แข็ง) ป้องกันหัวและคอ อยู่ห่างหน้าต่าง ผนังนอก เฟอร์นิเจอร์สูง อย่าวิ่งลงบันได/กระโดดจากตึกอย่างไร้สติ อยู่ในลิฟต์ให้ออกชั้นที่ใกล้ที่สุดทันที กลางแจ้งไปที่โล่ง อยู่ห่างสายไฟ ป้ายโฆษณา บ้านทรุด หลังแผ่นดินไหวตรวจแก๊สและไฟฟ้า ระวังอาฟเตอร์ช็อก
  \item อย่าปิดเสียงอย่าปิด: เตือนแผ่นดินไหวควรอนุญาตให้ "ทะลุผ่านโหมดเงียบ/ห้ามรบกวน" ดัง ก่อนนอนวางมือถือให้เอื้อมถึง อย่าหรี่เสียงเรียกเข้าเป็น 0
\end{itemize}

\booksubsection{เตือนแผ่นดินไหวในตัวของมือถือ: แต่ละแบรนด์เปิดอย่างไร}

ทางเข้าด้านล่างอาจต่างกันเล็กน้อยตามเวอร์ชันระบบ ค้นหา “แผ่นดินไหว” “เตือนภัย” “เตือนความปลอดภัย” ในการตั้งค่าได้โดยตรง ระบบ iOS รองรับการเตือนแผ่นดินไหวในจีนอย่างจำกัด มักต้องติดตั้งแอปเตือนภัยบุคคลที่สามที่เชื่อถือได้ และอนุญาตสิทธิ์ตำแหน่งและการแจ้งเตือน

\begin{center}
\begin{tabularx}{\linewidth}{|X|X|X|}
\hline
\textbf{แพลตฟอร์ม} & \textbf{ทางเข้าทั่วไป} & \textbf{คำอธิบาย} \tabularnewline
\hline
Xiaomi / Redmi & Mobile Manager / ศูนย์ความปลอดภัย > เตือนแผ่นดินไหว หรือ ค้นหา "แผ่นดินไหว" ในการตั้งค่า & บางเวอร์ชันมีบริการเตือนภัยในระบบ ต้องต่อเน็ตและเปิดตำแหน่ง \tabularnewline
Huawei / Honor & การตั้งค่า > ความปลอดภัย > แจ้งเตือนเหตุฉุกเฉิน / เตือนแผ่นดินไหว & ชื่ออาจเป็น "แจ้งเตือนเหตุฉุกเฉิน" ครอบคลุมแผ่นดินไหวด้วย \tabularnewline
OPPO / realme / OnePlus & การตั้งค่า > ความปลอดภัย (SOS ติดต่อฉุกเฉิน) > เตือนแผ่นดินไหว & บางรุ่นขึ้นกับบริการในพื้นที่ \tabularnewline
vivo / iQOO & แอปสภาพอากาศ หรือ การตั้งค่า > ความปลอดภัย/อื่น ๆ > เตือนแผ่นดินไหว & ยึดเมนูของเวอร์ชันระบบจริง \tabularnewline
iPhone / iOS & ในการตั้งค่าระบบมักไม่มีสวิตช์ในตัวครบถ้วน ติดตั้งแอปเตือนแผ่นดินไหวที่เชื่อถือได้และอนุญาตการแจ้งเตือน & แอปต้องทำงานเบื้องหลังและมีสิทธิ์ตำแหน่ง อย่าให้โหมดห้ามรบกวนของระบบบล็อกการเตือน \tabularnewline
Samsung เป็นต้น & การตั้งค่า > ความปลอดภัยและเหตุฉุกเฉิน / การเตือนฉุกเฉินไร้สาย & "การเตือนฉุกเฉิน" ครอบคลุมต่างกันตามพื้นที่และผู้ให้บริการ \tabularnewline
\hline
\end{tabularx}
\end{center}

\begin{tcbwarn}{ความเข้าใจผิดทั่วไปเรื่องเตือนแผ่นดินไหว}
\begin{itemize}
  \item เตือน $\neq$ พยากรณ์: บอกไม่ได้ว่า "พรุ่งนี้แผ่นดินไหวไหม" แจ้งเตือนเฉพาะเมื่อแผ่นดินไหวเกิดแล้วและคลื่นยังมาไม่ถึง
  \item ไม่ใช่ทุกที่จะได้ล่วงหน้าหลายวินาที: ใกล้ศูนย์กลางอาจมีแค่ไม่กี่วินาทีหรือไม่ทัน การเตือนหลักประโยชน์อยู่ที่พื้นที่รอบ ๆ
  \item ปิดการแจ้งเตือนเท่ากับไม่ได้ติดตั้ง: เตือนภัยดีแค่ไหน ถ้าเงียบ ปิดเบื้องหลัง ไม่ให้ตำแหน่ง ก็อาจไม่ได้รับ
  \item แอปปลอม: ระวังซอฟต์แวร์ "พยากรณ์แผ่นดินไหว" ที่ขอสิทธิ์รายชื่อและข้อความ ใช้ของระบบในตัวหรือของหน่วยงานที่เชื่อถือได้ก่อน
\end{itemize}

\end{tcbwarn}

\begin{tcbinfo}{จับคู่ความสามารถเตือนภัยอย่างไร}
แนะนำประจำวัน: “การแจ้งเตือนเตือนสภาพอากาศของระบบ + ช่องทางอุตุฯ/ฉุกเฉินทางการ + เตือนแผ่นดินไหว (ระบบหรือแอปที่เชื่อถือได้)” เหมือนการป้องกันสายรบกวน ไม่ต้องติดตั้งซอฟต์แวร์ประเภทเดียวกันหลายตัว กุญแจสำคัญคือเปิดสิทธิ์การแจ้งเตือน อนุญาตตำแหน่ง และทะลุโหมดห้ามรบกวนได้ หลังเปิดบล็อกสายรบกวนให้มือถือผู้ใหญ่แล้ว ตรวจเตือนสภาพอากาศและเตือนแผ่นดินไหวด้วย

\end{tcbinfo}

\begin{tcbpractice}{ลงมือทำ: ตรวจการตั้งค่าเตือนภัยใน 10 นาที}
\begin{enumerate}
  \item เปิดแอปสภาพอากาศของระบบ เข้าการตั้งค่า เปิดการแจ้งเตือน "เตือนสภาพอากาศ / สภาพอากาศภัยพิบัติ"
  \item ค้นหา "แผ่นดินไหว" หรือ "เตือนภัย" ในการตั้งค่าระบบ หาและเปิดเตือนแผ่นดินไหว/เตือนฉุกเฉิน (Android) สำหรับ iPhone ติดตั้งแอปเตือนแผ่นดินไหวที่เชื่อถือได้และอนุญาตการแจ้งเตือน
  \item ติดตามบัญชีทางการของสำนักอุตุฯ หรือหน่วยฉุกเฉินในพื้นที่ อนุญาตให้ส่งการแจ้งเตือน
  \item ยืนยันว่าเสียงเตือนไม่ถูกโหมดห้ามรบกวนเงียบสนิท ตรวจให้มือถือพ่อแม่ด้วย
\end{enumerate}

\end{tcbpractice}

\booksection{ความร่วมมือข้ามอุปกรณ์: ควบคุมคอมพิวเตอร์ ส่งเร็ว และความปลอดภัยบัญชี}

มือถือใช้งานราบรื่น ส่วนใหญ่เพราะมัน “จับคู่” กับคอมพิวเตอร์และแท็บเล็ต ส่วนนี้เสริมเครื่องมือใช้งานจริงหลายประเภท: ควบคุมคอมพิวเตอร์จากมือถือ AI ออฟฟิศและความร่วมมือระยะไกล วอลเปเปอร์เคลื่อนไหว ส่งไฟล์เร็วใน LAN คลิปบอร์ดข้ามอุปกรณ์ รหัสผ่านแบบไดนามิก 2FA และการเชื่อมโยงการสแกนกระดาษกับงานในตาราง การเชื่อมโยงหลายอุปกรณ์อย่างเป็นระบบดูบทที่ 17

\booksubsection{ควบคุมคอมพิวเตอร์จากมือถือ: Sunlogin / ToDesk}

\begin{itemize}
  \item Sunlogin Remote Control เหมาะกับ: ผู้ใช้ที่ต้องการควบคุม Windows หรือ Mac ที่บ้าน/บริษัทขณะอยู่ข้างนอก ใช้ช่วยผู้ใหญ่ซ่อมมือถือระยะไกลบ่อย (ดูบทที่ 27) จุดเด่น: ระบบนิเวศจีนเจริญแล้ว แอปมือถือควบคุมเดสก์ท็อปคอมพิวเตอร์ ส่งไฟล์ เปิดปิดเครื่อง (บางส่วนต้องใช้ฮาร์ดแวร์) เวอร์ชันฟรีเพียงพอสำหรับใช้ส่วนตัวเบา ๆ
  \item ToDesk เหมาะกับ: ผู้ใช้ที่ต้องการการเชื่อมต่อเสถียร หน่วงเวลาต่ำ ทำงาน/ดูแลระบบระยะไกล จุดเด่น: เชื่อมต่อด้วยรหัสอุปกรณ์+รหัสผ่าน ควบคุมคอมจากมือถือ ควบคุมมือถือจากคอม (แล้วแต่เวอร์ชัน) ส่งไฟล์และสลับหลายจอค่อนข้างสะดวก โหนดในจีนมักประสบการณ์ดี
  \item วิธีใช้พื้นฐาน: ติดตั้งไคลเอนต์บนคอมพิวเตอร์และเข้าสู่ระบบ/ออนไลน์ไว้ ติดตั้งแอปบนมือถือ ใช้รหัสอุปกรณ์ + รหัสผ่านเชื่อมต่อ (หรือบัญชีเดียวกัน) สร้างเซสชัน ใช้แทร็กแพด/ท่าทางบนมือถือควบคุมเดสก์ท็อป
  \item เส้นแดงความปลอดภัย (ต้องจำ): คนแปลกหน้าอ้างว่าเป็น "ฝ่ายบริการลูกค้า/ตำรวจ-อัยการ-ศาล/ผู้ให้บริการเครือข่าย" ให้คุณติดตั้ง Sunlogin, ToDesk และควบคุมมือถือหรือคอมพิวเตอร์ของคุณระยะไกล ปฏิเสธทั้งหมด—นี่คือวิธีหลอกลวงที่เป็นแบบแผน (ดูรายละเอียดในบทที่ 27)
  \item เมื่อใช้เอง: เปิดการยืนยันสองชั้น เปลี่ยนรหัสผ่านเริ่มต้น เมื่อไม่ใช้ให้ออกจากบัญชีหรือปิดโหมดไม่มีคนดูแล คอมพิวเตอร์บริษัทปฏิบัติตามกฎความปลอดภัยขององค์กร
  \item ทำงานระยะไกลในที่สาธารณะระวังคนแอบดูหน้าจอและรหัสผ่าน
\end{itemize}

\booksubsection{ซอฟต์แวร์ AI ออฟฟิศ: Trae / WorkBuddy / Dumate}

\begin{itemize}
  \item Trae เหมาะกับ: ผู้ใช้โปรแกรมมิ่งและเอกสาร ภาษาจีนเป็นมิตร จุดเด่น: AI IDE/ผู้ช่วยโปรแกรมมิ่งจาก ByteDance โหมด SOLO สร้างโปรเจกต์จากประโยคเดียว เขียนโค้ด แก้เอกสาร ประจำวันมีประสิทธิภาพ หากทำงานบนคอม ดู/ติดตามงานบนมือถือผ่านการควบคุมระยะไกลหรือความสามารถเสริมทางการได้ (ยึดรูปแบบผลิตภัณฑ์ล่าสุด) ดูรายละเอียดในบทที่ 41 ภาพรวมเครื่องมือ AI โปรแกรมมิ่ง
  \item WorkBuddy เหมาะกับ: คนทำงานที่อยากให้ AI ช่วยจัดบันทึกการประชุม อีเมล รายงาน เอกสารสำนักงาน จุดเด่น: ผู้ช่วย AI สำหรับออฟฟิศ เหมาะกับ "ร่าง—ขัดเกลา—สรุป" งานเขียนที่ทำบ่อย บางเวิร์กโฟลว์เริ่มบนมือถือได้ งานหนักให้ทำบนคอม
  \item Dumate เหมาะกับ: ทีม/บุคคลที่ต้องการความร่วมมือหลายคน จัดการเอกสารและความรู้ครบวงจร จุดเด่น: เครื่องมือออฟฟิศ AI ที่เน้นการรวบรวมข้อมูลและแก้ไขร่วมกัน ดูและแก้ไขเบา ๆ บนมือถือได้ การจัดวางซับซ้อนแนะนำทำบนจอใหญ่ของคอม
  \item "ควบคุมคอมพิวเตอร์จากมือถือ" เข้าใจอย่างไร: ความสามารถ AI ออฟฟิศจำนวนมากทำงานบนคอม/เว็บ เชื่อมกลับคอมผ่าน Sunlogin/ToDesk บนมือถือ หรือใช้แอป/เว็บของผลิตภัณฑ์เอง ก็ "บังคับ" งาน AI บนคอมขณะเดินทางไปทำงานได้ หลักการ: งานคำนวณหนัก เอกสารยาว วางคอม ดูความคืบหน้า แก้ไม่กี่บรรทัด วางมือถือ
  \item เตือนความเป็นส่วนตัว: ก่อนอัปโหลดสัญญา บัตรประชาชน ข้อมูลลูกค้า ดูนโยบายความเป็นส่วนตัวของผลิตภัณฑ์และข้อกำหนดการปฏิบัติตามขององค์กร ข้อมูลบริษัทใช้เครื่องมือที่องค์กรกำหนดก่อน
\end{itemize}

\booksubsection{วอลเปเปอร์เคลื่อนไหว: miHoYo Live Wallpaper (มือถือ + PC)}

\begin{itemize}
  \item Live Wallpaper เหมาะกับ: ผู้ใช้ที่ชอบตัวละคร miHoYo อยากได้วอลเปเปอร์เคลื่อนไหวฟรี จุดเด่น: แอปวอลเปเปอร์เคลื่อนไหวฟรีจาก miHoYo รองรับมือถือและ PC ตัวละคร Live2D/การแสดงเคลื่อนไหวใช้เป็นเดสก์ท็อปได้ บางรุ่น/ระบบตั้งเป็นหน้าจอล็อกหรือวอลเปเปอร์เดสก์ท็อปได้
  \item ประหยัดแบตและประสิทธิภาพ: วอลเปเปอร์เคลื่อนไหวเปลืองแบตกว่าภาพนิ่ง กิน GPU แบตต่ำหรือรุ่นเก่าใช้เฉพาะตอนชาร์จ หรือลดเฟรมเรต (ถ้าแอปมีตัวเลือก)
  \item แหล่งติดตั้ง: เว็บทางการ/แอปสโตร์ที่เชื่อถือได้ก่อน อย่า "เวอร์ชันดัดแปลง" จากบุคคลที่สามที่แถมโทรจัน บน PC ใส่ใจเส้นทางติดตั้งและการเริ่มต้นอัตโนมัติว่าจำเป็นไหม
\end{itemize}

\booksubsection{ส่งไฟล์เร็วใน LAN: LocalSend / IPMsg / AirDrop}

\begin{itemize}
  \item LocalSend เหมาะกับ: ผู้ใช้ Android/Windows/iOS/Mac/Linux ปนกัน อยากส่งข้ามแพลตฟอร์มโดยไม่ผ่านคลาวด์สาธารณะ จุดเด่น: โอเพ่นซอร์ส เชื่อมต่อตรงใน LAN ถ่ายโอนรูป วิดีโอ เอกสารใน Wi-Fi เดียวกัน ไม่ต้องใช้สาย ไม่ผ่านเซิร์ฟเวอร์ผู้ผลิต ความเร็วขึ้นกับเราเตอร์
  \item IPMsg เหมาะกับ: คอม Windows ในออฟฟิศ สถานการณ์ LAN เก่า จุดเด่น: เครื่องมือถ่ายโอน/แชทใน LAN เก่าแก่ เหมาะกับส่งไฟล์ใหญ่ในเครือข่ายภายใน ดาวน์โหลดจากช่องทางที่เชื่อถือได้ เครือข่ายภายในองค์กรปฏิบัติตามนโยบาย IT
  \item AirDrop เหมาะกับ: ผู้ใช้ iPhone, iPad, Mac จุดเด่น: แชร์ในระบบนิเวศ Apple ในคลิกเดียว ค้นพบด้วย Bluetooth + ถ่ายโอน Wi-Fi ตรง ไม่ใช้ข้อมูล ตั้ง "เฉพาะรายชื่อ" เพื่อลดการรบกวนจากคนแปลกหน้า
  \item อื่น ๆ: Mi Share, Huawei Share, vivo Share ฯลฯ เร็วมากในอุปกรณ์แบรนด์เดียวกัน "ผู้ช่วยส่งไฟล์" ของ WeChat สะดวกแต่ติดข้อจำกัดขนาดไฟล์และอุปกรณ์เข้าสู่ระบบ ไฟล์ใหญ่แนะนำ LocalSend/AirDrop
  \item เคล็ดลับการใช้: ทั้งสองฝ่ายต่อ Wi-Fi เดียวกัน ส่งไฟล์สำคัญบน Wi-Fi สาธารณะมีความเสี่ยง ใช้ฮอตสปอตส่วนตัวก่อน หลังถ่ายโอนเสร็จให้ฝั่งรับยืนยันว่าไฟล์ครบถ้วน
\end{itemize}

\booksubsection{ซิงก์คลิปบอร์ดข้ามอุปกรณ์: KDE Connect / ถ่ายโอนมือถือคอม}

\begin{itemize}
  \item Apple "Universal Clipboard" เหมาะกับ: ผู้ใช้ iPhone + Mac จุดเด่น: เข้าสู่ระบบ Apple ID เดียวกัน เปิด Handoff แล้ว คัดลอกบนมือถือ วางบน Mac ได้โดยตรง (กลับกันก็ได้) ข้อความและไฟล์บางอย่างใช้ได้
  \item KDE Connect เหมาะกับ: Android + Windows/Linux/Mac ผู้ใช้ที่อยากซิงก์คลิปบอร์ดโอเพ่นซอร์ส กระจกการแจ้งเตือน ป้อนข้อมูลระยะไกล จุดเด่น: จับคู่ใน LAN แล้วซิงก์คลิปบอร์ด ดันการแจ้งเตือนมือถือไปคอม ใช้มือถือเป็นรีโมทเปลี่ยนสไลด์/แทร็กแพด ข้อมูลไม่ผ่านบัญชีบุคคลที่สาม
  \item "ถ่ายโอนมือถือคอม" ของแบรนด์: Huawei Multi-screen Collaboration/Super Device, Xiaomi MiMo+, vivo Share, OPPO Connect, Honor Connect ฯลฯ มักมีคลิปบอร์ดข้ามหน้าจอและลากวางไฟล์ในตัว ประสบการณ์แบรนด์เดียวกันดีที่สุด
  \item ช่องทางทางการ Windows + Android: "Phone Link / Link to Windows" (พัฒนามาจาก Your Phone) ดูการแจ้งเตือนมือถือ รูปภาพบนคอม และคัดลอกวางในสถานการณ์ที่รองรับ Samsung ฯลฯ ผสานลึกกว่า
  \item เตือนความปลอดภัย: คลิปบอร์ดอาจมีรหัสผ่าน รหัสยืนยันค้างอยู่ บนคอมที่ใช้ร่วมกันระวังการซิงก์ทั่วไป ใช้เสร็จอย่าวางซ้ำในช่องรหัสผ่าน บางแอปแป้นพิมพ์ก็อ่านคลิปบอร์ด ใส่ใจสิทธิ์
\end{itemize}

\booksubsection{รหัสผ่านแบบไดนามิก 2FA: Aegis / Microsoft Authenticator}

\begin{itemize}
  \item ทำไมใช้ตัวยืนยันตัวตนแทนข้อความ: รหัสยืนยันทางข้อความอาจถูกจี้ SIM รหัสผ่านแบบไดนามิก TOTP สร้างในเครื่องตามอัลกอริทึมเวลา ไม่ต้องพึ่งสัญญาณเบอร์มือถือ ต้านฟิชชิงและการดักจับได้ดีกว่า (ยังต้องระวังโทรจันและภาพหน้าจอ)
  \item Aegis Authenticator เหมาะกับ: ผู้ใช้ Android ชอบโอเพ่นซอร์ส ควบคุมในเครื่อง จุดเด่น: โอเพ่นซอร์สฟรี รองรับสำรองข้อมูลเข้ารหัสและนำเข้าส่งออก อินเทอร์เฟซสะอาด ปลดล็อกด้วยลายนิ้วมือ/ใบหน้าสำหรับรายการโทเค็น
  \item Microsoft Authenticator เหมาะกับ: ผู้ใช้บัญชี Microsoft, Office 365/บัญชีองค์กร ข้าม iOS/Android จุดเด่น: อนุมัติการเข้าสู่ระบบแบบพุช ตรวจสอบรหัสผ่านและสำรองบัญชี (ตามเวอร์ชันผลิตภัณฑ์) เชื่อมโยงกับ Windows/ระบบตัวตนองค์กรอย่างแน่นแฟ้น
  \item Google Authenticator / 2FAS และอื่น ๆ: เจอบ่อยในกลุ่มนี้เหมือนกัน สิ่งสำคัญคือต้องมีการสำรองข้อมูลแบบเข้ารหัส เวลาเปลี่ยนเครื่องจะได้ไม่เสียโค้ดยืนยัน (token)
  \item สิ่งที่ต้องทำเมื่อเปลี่ยนเครื่อง: ย้ายข้อมูลจากเครื่องเก่า หรือใช้ "รหัสสำรอง/รหัสกู้คืน" ในบริการนั้น พิมพ์รหัสสำรองเก็บไว้สองที่ อย่าเก็บไว้ในมือถือเครื่องเดียว
  \item ส่วนเรื่องตัวจัดการรหัสผ่าน: คลังรหัสผ่านบางตัวก็เก็บ TOTP ได้ สะดวกดีแต่ความเสี่ยงจะกระจุกตัวมากขึ้น บัญชีที่สำคัญแนะนำให้แบ่งหน้าที่กันด้วย "ตัวจัดการรหัสผ่าน + แอปยืนยันตัวตนแยกต่างหาก" ดูรายละเอียดใน บทที่ 14
\end{itemize}

\booksubsection{คอมโบที่แนะนำ: วิธีต่อกันให้เป็นหนึ่งวัน}

\textbf{ตัวอย่างเวิร์กโฟลว์ที่ใช้งานได้จริง}\par

\begin{itemize}
  \item ตอนเช้า: เปิด TickTick ดูตารางวันนี้ แวะดูสภาพอากาศ/การแจ้งเตือน เปิดคอมทำงานกับ Trae/เอกสาร แล้ววางมือถือไว้ข้าง ๆ
  \item ตอนกลางวัน: เอกสารกระดาษใช้แอปสแกนแปลงเป็น PDF แล้วส่งเข้าคอมผ่าน LocalSend/AirDrop ส่วนงานเขียนยาว ๆ ร่างบนคอมด้วยเครื่องมือ AI สำหรับงานสำนักงาน
  \item ตอนออกไปข้างนอก: ใช้ ToDesk/Sunflower เชื่อมกลับไปที่คอมเพื่อทำงานต่อ ส่วนการล็อกอินบัญชีสำคัญ ๆ ใช้รหัสยืนยันแบบเปลี่ยนตลอดเวลา (OTP) จาก Aegis/Microsoft Authenticator
  \item ตอนกลางคืน: กลับมาเช็กถูกรายการที่ทำเสร็จใน TickTick และถ้าอยากผ่อนคลาย ก็เปลี่ยนวอลเปเปอร์หน้าจอให้ตัวเองสักหน่อย
\end{itemize}

\begin{tcbwarn}{คำจำง่ายเรื่องการควบคุมเครื่องระยะไกลและความปลอดภัยบัญชี}
\begin{itemize}
  \item ถ้าใครชวนให้ติดตั้งโปรแกรมควบคุมเครื่องระยะไกล วางสายทันที แล้วหาฝ่ายบริการลูกค้าทางการเอง
  \item ไม่บอกรหัสยืนยันหรือ OTP ให้ใครเด็ดขาด รวมถึงคนที่อ้างเป็น "พนักงานขนส่ง" หรือ "ฝ่ายบริการลูกค้า" ด้วย
  \item ก่อนรีโมตเข้าคอมบริษัทหรือติดตั้งซอฟต์แวร์ AI ให้ถามฝ่าย IT ก่อนว่าทำได้หรือไม่
  \item ระวังเป็นพิเศษเวลาส่งไฟล์หรือล็อกอินระยะไกลบนเครือข่ายสาธารณะ
\end{itemize}

\end{tcbwarn}

\begin{tcbpractice}{แบบฝึกหัดลงมือทำ: ลองเดินสายงาน "ข้ามอุปกรณ์" ให้ผ่านครบรอบหนึ่งครั้ง}
\begin{enumerate}
  \item ต่อมือถือกับคอมเข้า Wi-Fi เดียวกัน แล้วใช้ LocalSend (หรือ AirDrop) ส่งรูปหากันไปมาหนึ่งรูป
  \item สร้างงานใน TickTick ว่า "คืนนี้ 20:00 เก็บระเบียบไฟล์ที่สแกนแล้ว" แล้วยืนยันว่าทั้งมือถือและคอมเห็นงานนี้
  \item เปิด 2FA ให้อีเมลหรือ WeChat (ถ้ารองรับ) โดยผูกกับ Aegis หรือ Microsoft Authenticator แล้วเก็บรหัสสำรองไว้ด้วย
  \item ติดตั้ง Sunflower หรือ ToDesk ลงบนคอมของตัวเอง ลองเชื่อมต่อด้วยมือถือตัวเองให้สำเร็จหนึ่งครั้ง จากนั้นออกจากโปรแกรมหรือเสริมความแน่นหนาให้รหัสผ่าน เพื่อฝึกให้คุ้นมือโดยไม่ทิ้งช่องโหว่ความปลอดภัยไว้
\end{enumerate}

\end{tcbpractice}
